\chapter{Revisão da Literatura}

Esta seção tem como objetivo abordar os fundamentos teóricos que 
sustentam a pesquisa, divididos entre a capacidade técnica das ferramentas 
generativas, o comportamento humano perante a automação e o impacto destas 
tecnologias no fluxo de desenvolvimento. 

\section{A capacidade técnica e os limites da IA na geração de código} 

A introdução de Modelos de Linguagem de Grande Porte (LLMs) treinados 
especificamente em vastos repositórios de código-fonte representou um marco na 
automação do desenvolvimento de software. O estudo sobre o OpenAI Codex 
\cite{chen2021} demonstrou que essas ferramentas possuem uma notável 
eficiência na tradução rápida de especificações em linguagem natural para código. 
No entanto, é fundamental distinguir a agilidade na geração da real eficácia 
estrutural do código, isto é, a garantia de sua corretude e segurança funcional. Em 
cenários de escopo isolado, como a criação de funções algorítmicas, o 
desenvolvimento de scripts únicos ou a estruturação inicial de componentes de 
interface, a inteligência artificial atinge altas taxas de acerto. Essa capacidade de 
atuar como um autocompletar avançado acelerou a resolução de tarefas rotineiras.

Contudo, a precisão observada em testes de laboratório limitados não se 
traduz linearmente para ambientes de produção corporativos. A literatura atual 
destaca uma distinção fundamental entre a correção sintática e a validade semântica 
do código gerado. Modelos generativos frequentemente produzem blocos de código 
que parecem corretos e muitas vezes compilam sem erros, mas que falham na 
execução ou na lógica de negócio. Essa ineficiência ocorre através da invocação de 
métodos inexistentes em determinadas versões de frameworks, importações de 
bibliotecas depreciadas ou injeção de lógicas que desrespeitam o fluxo arquitetural 
\cite{liu2023}. Essa limitação está fortemente atrelada à restrição de contexto 
dos modelos, que ainda falham em manter a coerência quando a solução exige o 
mapeamento simultâneo de dezenas de dependências espalhadas pelo projeto.

O contraste entre o desempenho em escopos isolados e a engenharia de 
software do mundo real se destaca em avaliações rigorosas, como o SWE-bench 
\cite{jimenez2024}. Este teste avalia a capacidade da IA de resolver erros e 
melhorias reais reportados na plataforma GitHub\footnote{O GitHub é uma plataforma colaborativa de hospedagem de código-fonte e controle de versão baseada no sistema Git, amplamente utilizada globalmente para o gerenciamento de repositórios e o desenvolvimento de software em equipe.}. Os resultados revelam que, ao 
serem desafiadas a modificar bases de código legadas e interligadas, onde uma 
alteração em um banco de dados relacional pode impactar diversas rotas de uma 
API ou outras áreas de uma infraestrutura, as taxas de sucesso das IAs despencam 
consideravelmente. A resolução de falhas estruturais exige não apenas a escrita de 
novas linhas, mas a compreensão profunda do ecossistema do projeto, algo que as 
ferramentas atuais não conseguem abstrair isoladamente. 

Assim, embora as IAs generativas sejam ótimos assistentes, elas 
estão longe de funcionar de maneira isolada e autônoma em arquiteturas de 
sistemas complexos \cite{jimenez2024}. É indispensável a análise profunda e 
a aplicação de testes para impedir que soluções estruturalmente falhas ou 
vulnerabilidades sejam integradas ao ambiente de produção \cite{pearce2022, perry2023}.

\section{O fator humano, viés de automação e segurança}
A transição do papel do engenheiro de software, que gradativamente passa 
de autor primário do código para revisor de soluções geradas por algoritmos, traz 
implicações cognitivas profundas. A psicologia cognitiva descreve o cérebro humano 
como um avarento cognitivo, um sistema que busca constantemente otimizar o esforço mental por meio de atalhos \cite{fiske1984}. Quando uma 
ferramenta baseada em Inteligência Artificial fornece uma resposta estruturada e 
sintaticamente atraente, ocorre o fenômeno conhecido como viés de automação. 
Mosier e Skitka (1998) definem este viés como a tendência do indivíduo de utilizar a 
sugestão de um sistema automatizado como uma substituição para a busca ativa e 
vigilante de informações. Na prática, o desenvolvedor aceita a solução da máquina 
como correta por padrão, diminuindo o rigor do questionamento e ignorando falhas 
que seriam evidentes caso ele próprio estivesse a arquitetar o código desde o início.  

Esse relaxamento do pensamento crítico agrava-se durante a etapa de 
revisão de código. Revisar o código escrito por outro humano envolve compreender 
a lógica, a intenção e o contexto de um raciocínio construído. Por outro lado, revisar 
blocos maciços de código gerados instantaneamente por um LLM exige um nível 
exaustivo de concentração para identificar lógicas invertidas ou complexas demais, 
as chamadas “alucinações”, momentos em que a IA começa a gerar erros 
contextuais \cite{liu2023}, ou falhas de integração com a arquitetura existente do 
repositório \cite{jimenez2024}. Com o cansaço mental natural do processo de 
desenvolvimento, a capacidade de o programador manter a atenção cai, resultando 
em uma descarga cognitiva que afeta o pensamento crítico \cite{gerlich2025} e a 
probabilidade de aprovar códigos que parecem funcionar superficialmente, mas 
quebram em certos cenários específicos aumenta exponencialmente.

As consequências deste comportamento para a segurança da informação 
são alarmantes. Perry et al. (2023) investigaram a ilusão de segurança na adoção de assistentes de IA e 
demonstraram em experimentos controlados que programadores com acesso a ferramentas generativas não apenas 
tendem a produzir códigos com mais vulnerabilidades, mas também reportam uma 
confiança significativamente maior de que os seus códigos estão corretos e seguros. 

A literatura aponta que há também o risco da aversão ao uso da IA , e isto representa um desafio. 
Caso o desenvolvedor presencie falhas  graves ou sucessivas geradas pela IA, ele pode desenvolver uma 
aversão ao algoritmo, mesmo em momentos em que pode ser usada a seu favor de 
maneira saudável \cite{dietvorst2015}.

A intervenção humana diante das situações que ocorrem ao utilizar IA no 
desenvolvimento de software de maneira displicente se torna um padrão necessário 
no processo. A confiança excessiva e a descarga cognitiva comprometem 
diretamente a eficácia da auditoria manual \cite{gerlich2025,mosier1998}.
 Diminuir estes riscos exige que a avaliação qualitativa de código deixe de ser 
uma leitura passiva \cite{pearce2022}  e passe a integrar metodologias de 
validação rigorosas, partindo da premissa de que o código gerado por LLMs não é 
confiável até que a sua segurança e lógica sejam provadas \cite{perry2023,liu2023}.

\section{Impactos na produtividade e no mercado de trabalho}

A integração de agentes de inteligência artificial generativa alterou
substancialmente o trabalho do desenvolvedor de software. A transição da escrita
manual para geração automatizada possibilitou um ganho imediato no volume de
código produzido. Estudos empíricos demonstram que a adoção de ferramentas
como o GitHub Copilot aumenta a velocidade de conclusão de tarefas e
produtividade geral dos desenvolvedores \cite{peng2023}.

Contudo, o aumento da produção de código não se traduz necessariamente
em software seguro. Com a IA assumindo a codificação bruta quase que
instantaneamente, o gargalo produtivo do mercado deslocou-se irreversivelmente da
etapa de escrita para as etapas de revisão, integração e auditoria \cite{demirer2026}. Consequentemente, o mercado de trabalho passa a exigir
um novo perfil analítico. O profissional deixa de ser valorizado estritamente pela sua
capacidade de redigir algoritmos do zero e passa a ser exigido pela sua proficiência
em auditar lógicas complexas e mapear integrações de arquitetura, evitando a
aceitação passiva de código vulnerável.

Em suma, a inteligência artificial e o esforço humano apresentam, segundo
as análises de mercado, uma relação de forte complementaridade. Em vez de
eliminar postos de trabalho, a tendência corporativa é a busca por profissionais que
dominem a calibração dessas ferramentas e que atuem como validadores
essenciais, assegurando que o ganho de velocidade no desenvolvimento não resulta
na degradação da qualidade técnica e estrutural dos projetos a longo prazo \cite{demirer2026}.

\section{Trabalhos relacionados e comparação metodológica}

A avaliação da qualidade de código-fonte gerado por modelos de linguagem de grande escala tem sido objeto de diversas abordagens metodológicas na literatura recente. Estudos de base quantitativa, como o de \citeonline{liu2023}, fundamentam-se em \textit{benchmarks} automatizados, a exemplo do HumanEval, que mensuram a quão certeiros e corretos são os algoritmos em escopos isolados por meio de testes de unidade. 

Aumentando o nível de complexidade, pesquisas como a de \citeonline{jimenez2024} propõem metodologias baseadas na resolução de problemas em ecossistemas reais. Utilizando o \textit{SWE-bench}, os autores testam a capacidade dos LLMs de solucionar \textit{issues} de repositórios do GitHub, focando em métricas de taxa de resolução automatizada de problemas interligados em bases de código legadas. Já no campo da segurança da informação, \citeonline{perry2023} e \citeonline{pearce2022} utilizam simulações experimentais com desenvolvedores e análise de código para catalogar a injeção quantitativa de vulnerabilidades estruturais (como injeção de SQL e estouro de \textit{buffer}).

Embora a literatura existente forneça uma base sólida sobre a eficiência e as taxas de erro dessas ferramentas, grande parte dessas metodologias concentra-se em validações automatizadas ou em métricas estritamente quantitativas. Em contraste, a proposta adotada neste trabalho diferencia-se ao aplicar uma avaliação qualitativa humana ancorada nos critérios da norma ISO/IEC 25010 (Adequação Funcional, Confiabilidade, Manutenibilidade e Segurança). 


A avaliação científica dos limites das ferramentas generativas de código-fonte desenvolveu-se recentemente sob diferentes metodologias, conforme sumarizado no Quadro 1.

\begin{quadro}[htb]
\centering
\caption{Comparativo de Metodologias de Avaliação de Geração de Código por IA.}
\label{quadro:comparativo}
\begin{tabular}{|p{2.8cm}|p{2.5cm}|p{2.8cm}|p{3.5cm}|p{3.1cm}|}
\hline
\textbf{Estudo / Trabalho} & \textbf{Tipo de Avaliação} & \textbf{Foco Principal} & \textbf{Vantagens Metodológicas} & \textbf{Limitações} \\ \hline
\textbf{HumanEval+ / EvalPlus} \cite{liu2023} & Automatizada por Testes de Unidade & Corretude funcional algorítmica de escopo fechado. & Escala massiva de geração de entradas de teste (fuzzing/mutação). Reduz pass@k em até 28,9\%. & Limita-se a problemas puramente matemáticos/lógicos isolados sem dependências externas. \\ \hline
\textbf{SWE-bench} \cite{jimenez2024} & Automatizada por Execução em Repositórios & Resolução de \textit{issues} complexas do mundo real no GitHub. & Alta fidelidade técnica em bases de código legadas gigantescas de projetos reais. & Custoso computacionalmente; taxas de resolução historicamente inferiores a 2\%. \\ \hline
\textbf{Asleep at the Keyboard} \cite{pearce2022} & Automatizada por Análise Estática (CodeQL) e Manual & Injeção de vulnerabilidades de segurança (MITRE CWE Top 25). & Mapeia de forma ampla 18 tipos de falhas de segurança em 54 cenários de prompt. & Estático; não captura a dinâmica interativa do programador humano com a IA. \\ \hline
\citeonline{perry2023} & Empírica com Usuários em Laboratório & Impacto do uso de assistentes de IA na segurança do código e na confiança do usuário. & Condução de ensaios comportamentais reais, evidenciando o viés de automação e excesso de confiança. & Escopo de problemas pequeno (apenas 5 tarefas analisadas manualmente). \\ \hline
\textbf{Este Trabalho (Proposta)} & Qualitativa Humana por Norma Técnica & Avaliação sistemática de qualidade multidimensional baseada na \textbf{ISO/IEC 25010}. & Integração direta da psicologia cognitiva (viés de automação) com critérios formais de engenharia de software. & Amostragem de código e escopo menor do que os testes totalmente automatizados. \\ \hline
\end{tabular}
\fonte{O autor (2026).}
\end{quadro}

Como demonstrado no Quadro 1, as abordagens existentes concentram-se fortemente em métricas puramente automatizadas de correção algorítmica \cite{liu2023} ou em análises estatísticas isoladas de vulnerabilidades estáticas \cite{pearce2022}. \citeonline{perry2023} aproximam-se da vertente humana, mas limitam-se a medir a ocorrência pontual de falhas de segurança.

A metodologia adotada neste TCC diferencia-se ao preencher uma lacuna metodológica crucial: a aplicação sistemática de uma avaliação holística baseada na norma ISO/IEC 25010 (desdobrada em Adequação Funcional, Confiabilidade, Manutenibilidade e Segurança). Em vez de quantificar friamente o percentual de acertos de testes automatizados, esta pesquisa analisa detalhadamente o estilo e a estrutura do código que os LLMs produzem, estabelecendo o paralelo imediato com os comportamentos cognitivos disfuncionais (como a descarga cognitiva apontada por \citeonline{gerlich2025} e a queda de engajamento cognitivo em testes randomizados controlados observada por \citeonline{georgiou2025}). O diferencial reside em demonstrar exatamente de que forma o viés de automação compromete cada uma das subcaracterísticas de qualidade do software final.


\subsection{Detalhamento Técnico das Abordagens Experimentais da Literatura}

Para compreender a fundo os limites das ferramentas generativas de código-fonte, faz-se necessário realizar um exame rigoroso dos arranjos e metodologias experimentais aplicados nos principais estudos de referência. Cada trabalho selecionado aborda a avaliação sob diferentes perspectivas de escala, automação e interação humana.

\subsubsection{Avaliação de Corretude Algorítmica e Robustez de Testes: EvalPlus \cite{liu2023}}

O estudo de \citeonline{liu2023} fundamenta-se no diagnóstico de que os benchmarks tradicionais de geração de código, como o \textit{HumanEval} da OpenAI, possuem uma grave limitação de suficiência de teste (apresentando uma média modesta de apenas 9,6 casos de teste por tarefa). Essa escassez quantitativa induz a um falso estado de confiança, fazendo com que códigos gerados por LLMs que contêm falhas em casos de borda sejam erroneamente classificados como corretos.

Para solucionar esse problema e revelar a real capacidade de programação das ferramentas, os autores construíram o framework \textit{EvalPlus} para expandir o benchmark original para o \textit{HumanEval+}, elevando a escala de casos de teste em 80 vezes (alcançando uma média de 764,1 testes por tarefa). O delineamento experimental deste estudo utilizou um pipeline híbrido composto por duas etapas principais:
\begin{enumerate}
    \item \textbf{Geração de Seeds via LLM (ChatGPT):} O modelo foi instruído a realizar uma análise de caixa-branca sobre o código de gabarito de cada uma das 164 tarefas do benchmark, extraindo seus tipos, formatos e restrições semânticas para propor cerca de 30 entradas de teste altamente refinadas focadas em condições de erro.
    \item \textbf{Mutação de Entradas Baseada em Tipos:} A partir das sementes de alta qualidade geradas pelo ChatGPT, o framework executou uma rotina de testes mutativos recursivos sobre as estruturas de dados de entrada (como incremento/decremento de inteiros, manipulação de limites em cadeias de caracteres e alteração estrutural em dicionários e listas). Para evitar falhas lógicas e falsos-positivos na execução, foram aplicados contratos de programa para filtrar entradas inválidas antes de submetê-las ao teste.
\end{enumerate}

A validação final ocorreu por meio de testes diferenciais , comparando o comportamento de 26 modelos de linguagem comerciais e de código aberto frente às respostas de gabarito estruturadas. Para agilizar o processo de teste sem perder eficácia, os autores formularam o \textit{HumanEval+-Mini} aplicando um algoritmo de otimização de cobertura de conjunto  focado em três métricas: cobertura de ramificação, eliminação de mutantes lógicos inseridos artificialmente no gabarito e falha empírica de amostras de modelos. Esse processo destilou a suíte para uma média compacta de 16,1 testes por tarefa, preservando a capacidade de detecção de defeitos.

Como conclusão técnica do experimento, \citeonline{liu2023} identificaram que 18 gabaritos do próprio benchmark \textit{HumanEval} original (11\% das tarefas lógicas) continham erros crassos de implementação (lógica ruim, falta de tratamento de limites ou problemas graves de desempenho). Ao reavaliar os LLMs na suíte robustecida do \textit{HumanEval+}, o desempenho médio reportado  despencou de forma severa, sofrendo reduções de até 19,3\% para métricas de topo  e 28,9\% para volumes de amostragem.

\subsubsection{Resolução de Defeitos em Repositórios Complexos: SWE-bench \cite{jimenez2024}}

O trabalho de \citeonline{jimenez2024} parte da premissa de que a engenharia de software do mundo real raramente envolve a escrita de algoritmos curtos em arquivos isolados. Resolver problemas reais demanda a navegação por repositórios complexos, compreensão de dependências modulares entre múltiplos arquivos e geração de correções estruturadas (patches). Para testar essa habilidade, os autores introduziram o benchmark \textit{SWE-bench}, composto por 2.294 problemas de engenharia extraídos diretamente de problemas (\textit{issues}) reais e solicitações de pull (\textit{pull requests} - PRs) resolvidas em 12 repositórios populares de Python de código aberto (como Django, SymPy, Scikit-Learn e Matplotlib).

A metodologia de construção e validação dos dados baseou-se em um pipeline automatizado de três etapas:
\begin{enumerate}
    \item \textbf{Raspagem e Seleção:} Extração de aproximadamente 90.000 PRs mescladas das bases do GitHub, filtrando aquelas que resolveram de forma clara um ticket de bug e que alteraram arquivos de teste associados.
    \item \textbf{Delineamento do Harness de Execução Hermético:} Para cada problema, o framework isolou o ambiente em contêineres e instalou a base de código correspondente no commit de referência anterior à correção humana. 
    \item \textbf{Filtro por Execução Dinâmica:} Os testes de unidade do próprio repositório foram executados antes e depois da aplicação do patch histórico do desenvolvedor. A tarefa só foi incluída no dataset final se apresentou comportamento de transição dinâmico (\textit{fail-to-pass} - F2P), onde pelo menos um teste de unidade passou a funcionar após a inserção do código de correção, excluindo cenários com falhas de instalação ou dependências inacessíveis.
\end{enumerate}

No ensaio de avaliação, os modelos recebiam apenas a descrição em texto do problema e o acesso à árvore do projeto. Devido aos limites de janela de contexto das redes, o experimento testou abordagens baseadas em recuperação esparsa (BM25) para delimitar as linhas de arquivos de interesse em blocos de até 27.000 tokens, além de um cenário de controle idealizado no qual apenas os arquivos realmente editados na solução humana foram fornecidos. O modelo gerava um arquivo de patch unificado (\textit{.patch}) que era injetado e executado contra os testes originais da suíte e mais 51 testes adicionais de preservação funcional (\textit{pass-to-pass} - P2P).

Os resultados experimentais demonstraram que a atividade de engenharia autônoma é extremamente desafiadora: o modelo proprietário mais robusto avaliado (Claude 2) solucionou apenas 1.96\% das tarefas com busca de arquivos e 4,80\% no arranjo idealizado com arquivos pré-selecionados. Já o ChatGPT obteve apenas 0,52\% de resoluções. A análise técnica das falhas revelou que a maioria das gerações caía em cenários de regressão (\textit{No-Op} ou quebra de componentes adjacentes), evidenciando que os LLMs sofrem por não mapearem interdependências profundas de arquitetura de software de grande porte.

\subsubsection{Investigação Sistemática de Vulnerabilidades Estáticas: Asleep at the Keyboard \cite{pearce2022}}

Visando caracterizar o potencial de risco de segurança na adoção cega de ferramentas generativas de código baseadas em dados públicos sem curadoria, \citeonline{pearce2022} executaram o primeiro ensaio sistemático focado na identificação estática de falhas do catálogo MITRE CWE (\textit{Common Weakness Enumeration}) Top 25 em códigos gerados pelo GitHub Copilot (durante sua fase de \textit{technical preview}).

O arranjo experimental estruturou-se sob a criação de 89 cenários sintéticos de desenvolvimento (\textit{prompts} parciais incompletos) projetados de modo que uma conclusão meramente lógica ou funcional induzisse o assistente à injeção de vulnerabilidades graves. O experimento analisou um montante de 1.689 programas gerados ao longo de três eixos fundamentais:
\begin{enumerate}
    \item \textbf{Diversidade de Fraquezas (DOW):} Consistiu no desenvolvimento de 54 cenários focados em 18 falhas do ranking CWE Top 25 (como estouro de buffer CWE-787, Cross-Site Scripting CWE-79 e Path Traversal CWE-22) utilizando as linguagens de programação C e Python.
    \item \textbf{Diversidade de Prompts (DOP):} Avaliou o comportamento dinâmico do modelo frente a alterações cosméticas e estruturais na escrita do prompt de controle para injeção de SQL (CWE-89), incluindo variações de metadados como autoria de código, comentários de documentação e ordem de parâmetros de execução.
    \item \textbf{Diversidade de Domínio (DOD):} Investigou a extrapolação do assistente para domínios industriais menos recorrentes em sua base de dados de treinamento, instanciando 18 cenários lógicos voltados ao desenvolvimento de hardware em descrição RTL (\textit{Register-Transfer Level}) por meio da linguagem Verilog (focando em 6 CWEs específicos de hardware).
\end{enumerate}

Para a avaliação de segurança do material gerado, os autores empregaram a suíte comercial de análise estática CodeQL da GitHub configurada especificamente para identificar cada fraqueza sintática isolada nos arquivos de software compiláveis de C e scripts de Python, complementando o processo com auditoria de segurança estritamente manual conduzida pelos próprios autores para falhas arquiteturais complexas e para os cenários em Verilog.

Os dados empíricos demonstraram que, de modo geral, aproximadamente 40,00\% de todas as soluções construídas automaticamente pelo Copilot eram inseguras. O comportamento foi severo na linguagem C, onde 50,29\% dos arquivos continham vulnerabilidades de escrita e manipulação de memória e, mais preocupante, 52,00\% das sugestões indicadas pela ferramenta com maior taxa de confiança lúdica continham brechas críticas. O ensaio DOP apontou alta sensibilidade do modelo a fatores de indução de contexto: adicionar uma tag indicando que o autor do arquivo era Andrey Petrov (desenvolvedor de renome da comunidade Python) reduziu quantitativamente as falhas e aumentou o índice de segurança das respostas.

\subsubsection{Ensaios Comportamentais de Interação e Viés Cognitivo: Perry et al. \cite{perry2023}}

O trabalho de \citeonline{perry2023} destaca-se como a pesquisa pioneira no estudo comportamental da interação humana direta com ferramentas de assistência a código baseadas em modelos fundamentais de linguagem (no caso, utilizando o motor \textit{codex-davinci-002} da OpenAI). Em vez de avaliar a ferramenta de forma isolada, os cientistas projetaram um laboratório focado em entender se o acesso à IA realmente induz programadores à escrita de códigos menos seguros e se esse processo é catalisador de vieses cognitivos como o de automação.

A arquitetura metodológica do experimento estruturou-se sob as seguintes definições:
\begin{itemize}
    \item \textbf{Composição de Participantes:} Recrutamento de 47 participantes de espectros variados de competência, desde estudantes universitários de ciência da computação até engenheiros de software profissionais com atuação de mercado.
    \item \textbf{Desenvolvimento do Instrumento de Estudo (UI):} Os cientistas programaram um ambiente de desenvolvimento integrado (IDE) autônomo e customizado construído sob as bibliotecas React, Redux e Electron. A interface para o grupo experimental continha uma aba de assistência integrada na qual os participantes podiam interagir de forma flexível com o modelo de IA através de buscas em linguagem natural, além de modificar manualmente os hiperparâmetros de inferência como temperatura (diversidade lúdica da saída) e comprimento máximo do bloco gerado. O grupo de controle, por sua vez, resolveu as tarefas utilizando apenas uma janela do navegador de internet para consultas gerais.
    \item \textbf{Casos de Uso de Segurança Práticos:} Os desenvolvedores foram desafiados a codificar 5 rotinas de software sensíveis distribuídas em C, JavaScript e Python: (Q1) Cifragem simétrica de arquivos com biblioteca AES; (Q2) Assinatura eletrônica de mensagens com algoritmo ECDSA; (Q3) Tratamento de controle de acesso de diretório local restrito; (Q4) Execução de consultas SQL parametrizadas para evitar injeção de dados; (Q5) Codificação de rotina de conversão de tipos de variáveis numéricas inteiras para caracteres com ponteiros manuais em linguagem C.
\end{itemize}

O monitoramento do ensaio coletou de forma contínua dados detalhados de telemetria das consultas enviadas ao modelo Codex, o histórico de alterações incrementais nas caixas de texto, alterações de temperatura, gravações de áudio e captação de vídeo de tela das sessões de codificação. A avaliação de segurança e conformidade de cada código entregue foi conduzida de forma independente por dois avaliadores especializados, medindo a confiabilidade através da correlação Kappa de Cohen (que variou de 0,7 a 0,96 para corretude de lógica e 0,68 a 0,88 para segurança).

Os resultados expuseram de forma quantitativa que os desenvolvedores do grupo experimental que usaram a IA para escrever seus códigos geraram soluções significativamente menos seguras do que o grupo de controle em 4 das 5 tarefas avaliadas (Q1, Q3, Q4 e Q5). Paralelamente, os dados das pesquisas pós-teste revelaram a manifestação severa do viés de automação: os participantes auxiliados pela ferramenta apresentaram taxas de confiança estatisticamente superiores de que suas soluções estavam corretas e seguras, mesmo quando seus códigos continham vulnerabilidades graves (como ausência de vetores de inicialização aleatórios em rotinas de criptografia ou falhas no dimensionamento de buffers em ponteiros em C). O estudo constatou que aqueles programadores que não aceitaram as primeiras sugestões da IA e realizaram alterações incrementais e de engenharia de prompts (definindo assinaturas explícitas de funções auxiliares ou ajustando o parâmetro de temperatura para obter maior dispersão na amostragem funcional) obtiveram códigos com melhor desempenho e menor taxa de insegurança lógica.

\subsubsection{Métricas Temporais e Eficiência de Produção: Peng et al. \cite{peng2023}}

A pesquisa de \citeonline{peng2023} afasta-se da análise estrita de vulnerabilidade de código para focar na medição empírica do real ganho de produtividade produtiva (velocidade e conformidade mecânica) associado ao uso do GitHub Copilot no cotidiano de desenvolvedores profissionais de software, estruturando o primeiro ensaio clínico controlado randomizado do segmento.

A metodologia e o ecossistema do ensaio basearam-se nas seguintes frentes técnicas:
\begin{itemize}
    \item \textbf{Amostragem de Desenvolvedores:} Recrutamento controlado de 95 programadores freelancers ativos localizados na plataforma internacional de prestação de serviços Upwork. Todos assinaram contratos digitais rígidos de confidencialidade e conformidade de ensaios comportamentais estruturados.
    \item \textbf{Formulação de Grupo e Randomização:} Os participantes foram divididos de forma aleatória em um grupo de tratamento composto por 45 engenheiros de software (que receberam acesso de licença temporária ao GitHub Copilot acompanhado de um vídeo explicativo de 1 minuto detalhando o uso de complementação contextual) e um grupo de controle composto por 50 engenheiros (que desenvolveram o sistema sem acesso ao recurso gerador, permitindo o uso de documentação técnica, Stack Overflow e ferramentas normais de busca).
    \item \textbf{Tarefa de Software Controlada:} Os participantes foram instruídos a programar, no menor intervalo temporal possível, um servidor HTTP básico em JavaScript (Node.js). A tarefa exigiu a recepção de conexões locais, análise estrutural de cabeçalhos de requisição de arquivos, validação de códigos de retorno de estado (como retornos HTTP de erros 404, 400 ou 405) e manipulação estruturada do sistema de arquivos físico do computador.
    \item \textbf{Rastreabilidade Dinâmica via GitHub Classroom:} A tarefa foi disponibilizada por meio de ramificações individuais privadas criadas dinamicamente na plataforma educacional da GitHub. O sistema computava o instante exato de criação do repositório como marca temporal de início de execução. A integridade lógica das submissões era avaliada por uma suíte interna de 12 testes de validação funcionais inflexíveis e imutáveis inseridos no projeto. O intervalo total gasto pelo desenvolvedor foi indexado com precisão milimétrica a partir do momento de criação até a gravação e push da primeira submissão que passou com 100\% de aproveitamento nos 12 testes funcionais da pipeline.
\end{itemize}

Os resultados do experimento controlado demonstraram um efeito expressivo: o grupo de tratamento que desenvolveu o servidor HTTP sob o suporte direto do GitHub Copilot concluiu a tarefa com uma velocidade 55,8\% superior em comparação com o grupo de controle. Em termos de mediana de tempo, os desenvolvedores assistidos necessitaram de apenas 71,1 minutos para fechar o sistema sem erros lógicos nos testes, enquanto os programadores desprovidos da ferramenta precisaram de 160,9 minutos. O estudo também apontou efeitos de heterogeneidade: os ganhos relativos de tempo foram proporcionalmente maiores para programadores mais velhos, profissionais que possuem carga de dedicação de programação diária elevada e indivíduos com menor tempo bruto de experiência em programação, sugerindo que as IAs atuam como niveladoras lógicas de gargalos de codificação manual.



